\chapter{Synthèse décisionnelle}
\label{ch:synthese}

\begin{encadre}{Encadré 8.1 — Résultats du modèle et recommandations : ce qui relève de quoi}
Les réponses aux huit questions qui suivent s'appuient sur les résultats des parties I à III ; leur degré de fiabilité
est indiqué. Le programme proposé ensuite est une \textbf{recommandation} : il résulte d'un choix de pondération entre
objectifs (réduire le déficit sans augmenter la pauvreté ni sacrifier l'activité de long terme) qui n'est pas dicté par
le modèle. Un autre choix de pondération — priorité absolue à la dette (scénario D), à la croissance (B) ou à la
redistribution (C) — conduirait à d'autres programmes, également cohérents avec le modèle. Les choix de société
(niveau des services publics, partage de l'effort entre générations, entre revenus du travail et du patrimoine,
acceptabilité) relèvent du débat démocratique.
\end{encadre}

\section{Réponses aux huit questions}

\subsection*{1. Quels sont les principaux dysfonctionnements structurels révélés par le modèle ?}

\begin{enumerate}
\item \textbf{Un déséquilibre budgétaire qui s'auto-entretient.} Le déficit primaire (3,4 points de PIB en 2030 dans la
référence) et le refinancement de la dette à des taux proches de 5~\% font croître la dette plus vite que le PIB. La
charge d'intérêts double d'ici 2032 (de 2,5 à 5,1 points de PIB). \fiab{robuste}.
\item \textbf{Une croissance potentielle faible}, bornée par une productivité de 0,7~\% par tête (hypothèse COR) et une
population active qui plafonne vers 2035. \fiab{conventionnel} : c'est une hypothèse de la référence, cohérente avec les
tendances observées depuis 2010.
\item \textbf{Un vieillissement qui pèse sur la protection sociale} : 0,51 personne de 65 ans ou plus par personne de 20 à
64 ans en 2047 contre 0,39 en 2025 ; pensions de 13,5 à 15,7~\% du PIB. \fiab{robuste} (démographie).
\item \textbf{Une dépendance aux importations} (contenu importé de 27~\% de la consommation, toute l'énergie fossile) qui
rend le pouvoir d'achat et les finances publiques vulnérables aux chocs de prix : +50~\% sur l'énergie importée coûte
1,2~\% de PIB et 15 points de dette en dix ans. \fiab{plutôt robuste}.
\item \textbf{Des taux marginaux effectifs élevés entre 1 et 2 SMIC} (55 à 82~\%), qui réduisent fortement le gain d'une
hausse de salaire pour les ménages modestes en emploi. \fiab{robuste} (législation).
\item \textbf{Une désindustrialisation tendancielle} (industrie manufacturière de 9,6 à 8,4~\% du PIB en 2037) que les
instruments représentés dans le modèle infléchissent peu. \fiab{fragile}.
\end{enumerate}

\subsection*{2. Quelles sont les dépenses publiques les moins efficaces économiquement et socialement ?}

Le modèle \textbf{ne permet pas de répondre} à la question de l'efficacité intrinsèque des dépenses : il ne mesure ni la
qualité des services publics, ni leurs effets sur la santé, l'éducation, la sécurité ou la productivité de long terme,
et il ne distingue pas les fonctions (santé, éducation, défense, fonctionnement ont les mêmes effets). Il permet en
revanche de dire, pour un euro d'économie : (a) \textbf{quelles économies coûtent le moins en activité} : les transferts
monétaires non ciblés et la sous-indexation temporaire des pensions coûtent trois à sept fois moins d'activité par point
de déficit réduit que la réduction de la production publique (emplois, fonctionnement) ; (b) \textbf{quelles dépenses
ont l'effet le plus durable sur la production} : l'investissement public (élasticité 0,08) ; c'est donc la dépense la
plus coûteuse à réduire à long terme ; (c) \textbf{quelles dépenses réduisent le plus la pauvreté par euro} : les
prestations ciblées (prime d'activité, RSA, minimum vieillesse), alors que les transferts non ciblés ont un effet
dilué. \fiab{plutôt robuste} pour (a) et (c) ; \fiab{fragile} pour (b).

\subsection*{3. Quels prélèvements pénalisent le plus la production, l'emploi et l'investissement ?}

Par euro prélevé, à court terme (tableau~\ref{tab:rendement}) : les \textbf{cotisations employeurs}, la \textbf{taxe
carbone}, la \textbf{CSG} et l'\textbf{impôt sur le revenu des classes moyennes} sont les plus coûteux en emploi
(37\,000 à 48\,000 emplois pour 9 à 11\md{} prélevés) ; la \textbf{TVA} est moins coûteuse en emploi mais réduit le pouvoir d'achat de
tous, y compris des retraités et des allocataires ; l'\textbf{IS} est le prélèvement qui pèse le plus sur
l'investissement ($-$0,44~\% pour une hausse de 4,7 points), mais le modèle en sous-estime probablement le coût total
(localisation des profits et des sièges non représentée) ; les \textbf{impôts sur le patrimoine} (ISF, PFU,
successions) pèsent le moins sur l'activité dans le modèle, sous réserve des réactions des contribuables, dont l'ampleur
est très incertaine (rendement net de l'ISF de 3,6 à 6,2\md{} selon l'élasticité). Les \textbf{impôts de production}
paraissent sans coût dans le modèle, faute de représentation de leurs effets sur l'investissement et la localisation :
ce résultat ne doit pas être utilisé. \fiab{plutôt robuste} pour le classement des impôts sur les ménages ;
\fiab{fragile} pour les impôts sur les entreprises.

\subsection*{4. Quels mécanismes entretiennent le chômage et la pauvreté ?}

Dans le modèle : (a) \textbf{une demande insuffisante} à court terme (le chômage réagit à l'activité avec une élasticité
de l'emploi de 0,26 la première année) ; (b) \textbf{le coût réel du travail} à moyen terme (élasticité de l'emploi de
$-$0,4) ; (c) \textbf{la lenteur d'absorption des chocs d'offre de travail} : un report de l'âge ou un allongement de la
durée du travail augmente d'abord le chômage ; (d) pour la pauvreté, \textbf{le niveau des minima sociaux et le non-recours}
(un tiers des ayants droit au RSA et à la prime d'activité), et \textbf{la sensibilité des ménages modestes aux prix}
(énergie, TVA). Le modèle ne représente pas les mécanismes microéconomiques du chômage (appariement, qualifications,
mobilité, discriminations) ni la persistance de la pauvreté (trajectoires individuelles) : \fiab{non établi} au-delà
de ces canaux macroéconomiques.

\subsection*{5. Quels leviers permettraient d'augmenter durablement les revenus réels ?}

Durablement, seuls les leviers qui augmentent la \textbf{production par habitant} augmentent les revenus réels : la
productivité (que le modèle prend comme une hypothèse), l'investissement public productif (+0,12 à +0,35~\% de PIB à
vingt ans pour 10 à 30\md{} par an), l'emploi des seniors (+1,3~\% de PIB en 2047 pour un report à 65 ans) et un prix de
l'énergie importée plus bas ($-$20~\% : +0,5~\% de PIB et +1,3~\% de pouvoir d'achat). Les hausses de salaires ou de
prestations redistribuent le revenu sans l'augmenter durablement. \textbf{Point important} : dans les scénarios qui
reportent l'âge ou allongent la durée du travail, la production augmente mais le pouvoir d'achat \emph{par ménage}
n'augmente pas dans le modèle (B : $-$1,6~\% en 2037), car le gain prend la forme d'emplois supplémentaires et de prix plus
bas plutôt que de revenus plus élevés pour les ménages en place. \fiab{fragile}.

\subsection*{6. Quelles réformes permettraient de réduire simultanément le déficit et la dette ?}

Toute mesure qui améliore le solde réduit la dette dans le modèle, à une exception près : le report de l'âge, qui améliore
peu le solde et réduit les prix, ce qui augmente légèrement le ratio de dette (effet de dénominateur, \fiab{fragile}). Les
réformes qui réduisent le déficit au moindre coût en activité sont : les impôts sur le patrimoine, la TVA (avec
compensation ciblée pour les ménages modestes), une modération temporaire de l'indexation des prestations les plus
élevées, et les transferts non ciblés. \textbf{Mais aucune combinaison réaliste ne stabilise la dette si le taux d'intérêt
reste proche de 5~\%} (chapitre~\ref{ch:strategie}) : la stabilisation exige une amélioration du solde primaire de l'ordre
de 6 points de PIB, ou une baisse durable des taux, ou une croissance plus forte.

\subsection*{7. Quelles mesures permettraient de relancer la production nationale ?}

À court terme : l'investissement public (multiplicateur d'environ 0,6 la première année compte tenu du contenu importé),
les prestations aux ménages modestes, la baisse des cotisations employeurs (+0,36~\% de PIB pour 20\md, mais non
financée). À long terme : le report de l'âge de départ, l'investissement public et une énergie moins chère. La
réindustrialisation augmente la production (+0,45~\% en 2032 pour 10~\% d'importations relocalisées) \emph{si} elle est
obtenue sans hausse des prix de l'énergie ni protection commerciale ; le modèle ne dit pas comment l'obtenir. Les droits
de douane réduisent la production nationale (riposte, prix). \fiab{plutôt robuste} pour le court terme ; \fiab{fragile}
pour le long terme.

\subsection*{8. Quels arbitrages seraient nécessaires ?}

\begin{enumerate}
\item \textbf{Entre le rythme de réduction du déficit et l'activité} : chaque point de PIB de consolidation coûte de 0,3 à
0,7~\% de PIB la première année avec les instruments les moins récessifs ; une consolidation rapide de 3 points coûte une
récession (scénario D).
\item \textbf{Entre le pouvoir d'achat des ménages en place et la soutenabilité} : les instruments de consolidation les
moins récessifs (TVA, désindexation) sont ceux qui réduisent le pouvoir d'achat de tous ; les compenser pour les ménages
modestes coûte environ un quart de leur rendement.
\item \textbf{Entre générations} : reporter l'âge de départ ou modérer l'indexation fait porter l'effort sur les actifs âgés
et les retraités ; ne rien faire le reporte sur les générations futures par la dette.
\item \textbf{Entre revenus du travail et du patrimoine} : taxer davantage le patrimoine pèse peu sur l'activité dans le
modèle, mais son rendement est très incertain (réactions des contribuables).
\item \textbf{Entre souveraineté et coût} : produire en France ce qui est importé coûte plus cher ; le protéger par des
droits de douane fait payer les consommateurs et les exportateurs.
\item \textbf{Entre chômage présent et production future} : travailler plus (carrières, durée) augmente la production à
long terme mais d'abord le chômage.
\end{enumerate}

\section{Programme de réformes recommandé (scénario R)}

\paragraph{Principe de construction.} Le scénario R n'est pas issu directement de l'optimisation ; il en retient les
enseignements, avec quatre règles : (1) n'utiliser que des leviers dont l'effet a au moins un signe robuste ; (2) préférer
les instruments de consolidation les moins coûteux en activité ; (3) neutraliser l'effet sur la pauvreté par des
prestations ciblées ; (4) préserver ou renforcer les leviers d'offre de long terme (investissement public, emploi des
seniors). Il exclut les leviers dont le modèle ne voit pas le coût (hausse d'IS et d'impôts de production), ceux qui
reposent sur une hypothèse non estimée (baisse des allocations chômage), la durée hebdomadaire (coût en emploi à court
terme et hypothèse de compensation non estimée) et la protection commerciale (effet négatif). Il a été simulé avec le
modèle complet (fichier \code{rapport/resultats/verification\_R.json}).

\begin{table}[htbp]
\centering\scriptsize
\caption{Programme recommandé : chiffrage mesure par mesure (méta-modèle, effets isolés mis à l'échelle) et du programme
simulé ensemble (modèle complet). Solde en Md€ de 2025 : points de PIB $\times$ PIB 2025.}
\label{tab:programmeR}
\tabfit{../tableaux/programme_R}
\end{table}

\begin{small}
\begin{longtable}{@{}p{2.1cm}p{12.2cm}@{}}
\caption{Les réformes du programme recommandé, par ordre de priorité.}\label{tab:reformes}\\
\toprule
\endfirsthead
\toprule\endhead
\bottomrule\endfoot
\multicolumn{2}{@{}l}{\textbf{1. Protéger les ménages modestes (préalable aux mesures de consolidation)}} \\
Mesure & Prime d'activité +25~\%, RSA +25~\%, minimum vieillesse +25~\%, au 1\textsuperscript{er} janvier 2028. \\
Problème & Pauvreté à son plus haut niveau depuis 1996 (15,4~\% en 2024) ; les hausses de TVA prévues pèsent davantage sur les modestes. \\
Mécanisme & Revenu des déciles 1 à 6 $\uparrow$ ; propension à consommer élevée $\to$ soutien de l'activité. \\
Coût & 10,4\md{} (0,35 point de PIB en 2030). \\
Effets & Pauvreté $-$1,8 point (statique, mesures isolées) ; PIB +0,36~\% en 2030 ; emploi +48\,000 en 2037. \\
Risques & Non-recours ; effet sur l'offre de travail non représenté (RSA) ; prolongement de la zone de taux marginaux élevés (prime d'activité). \\
Robustesse & \fiab{plutôt robuste} (législation, microsimulation ; effets statiques). \\
\midrule
\multicolumn{2}{@{}l}{\textbf{2. Consolider par les prélèvements les moins récessifs}} \\
Mesure & TVA : +2 points au taux normal, +1 point aux taux réduits ; droits de succession +50~\% ; impôt sur la fortune (1~\% au-delà de 1,3\,M€, remplaçant l'IFI) ; PFU +15 points. Entrée en vigueur 2028. \\
Problème & Déficit structurel ; coût en activité des autres instruments. \\
Mécanisme & Prélèvements à large assiette (TVA) ou sur des ménages à faible propension à consommer (patrimoine) $\to$ effet limité sur la demande. \\
Rendement & 47\md{} statiques, 39\md{} nets des réactions des contribuables ; 1,13 point de PIB de solde en 2030, soit 34\md{} (mesures isolées). \\
Effets & PIB $-$0,35~\% en 2030, effet presque nul en 2037 ; prix +1,1~\% ; pauvreté +0,8 point (compensée par le volet 1). \\
Risques & Réactions des hauts patrimoines (rendement net de l'ISF et du PFU de $-$70~\% à +40~\% autour de la valeur centrale selon l'élasticité) ; acceptabilité ; effet inflationniste ponctuel. \\
Robustesse & \fiab{robuste} pour le solde ; \fiab{fragile} pour le rendement des impôts sur le patrimoine. \\
\midrule
\multicolumn{2}{@{}l}{\textbf{3. Maîtriser la dépense courante}} \\
Mesure & Fonctionnement $-$10\md{} par an (hors investissement) ; sous-indexation des pensions de 0,5 point par an. 2028-2032, à réexaminer ensuite. \\
Problème & Dépense publique en volume plus dynamique que l'activité ; pensions en hausse avec le vieillissement. \\
Mécanisme & Moindre dépense ; l'indexation réduit la masse des pensions de façon cumulative. \\
Rendement & 7,6\md{} en 2030 (0,25 point de PIB), davantage ensuite (indexation). \\
Effets & PIB $-$0,36~\% en 2030 ; pensions réelles $-$0,5~\% par an (environ $-$2,5~\% au bout de cinq ans). \\
Risques & Qualité des services publics (non mesurée) ; pouvoir d'achat des petites pensions (une protection des pensions les plus basses serait souhaitable ; non modélisée). \\
Robustesse & \fiab{robuste} pour le solde ; \fiab{non établi} pour les effets de long terme des dépenses réduites. \\
\midrule
\multicolumn{2}{@{}l}{\textbf{4. Investir}} \\
Mesure & Investissement public +10\md{} par an à partir de 2028 (infrastructures, rénovation, équipements). \\
Problème & Faible croissance potentielle ; seule dépense dont l'effet sur la production persiste dans le modèle. \\
Mécanisme & Demande à court terme ; capital public productif à long terme (élasticité 0,08). \\
Coût & 9,4\md{} en 2030 (0,31 point de PIB) ; dette +3,3 points en 2037. \\
Effets & PIB +0,20~\% en 2030, +0,17~\% en 2037, +0,12~\% en 2047. \\
Risques & Sélection des projets (élasticité nulle si mal choisis) ; contenu importé ; capacités du BTP. \\
Robustesse & \fiab{plutôt robuste} à court terme ; \fiab{fragile} à long terme (élasticité 0,04 à 0,12). \\
\midrule
\multicolumn{2}{@{}l}{\textbf{5. Allonger progressivement les carrières}} \\
Mesure & Âge légal porté à 65 ans, à raison de trois mois par génération à partir de 2028. \\
Problème & Vieillissement ; population active en recul après 2035. \\
Mécanisme & Plus d'actifs et d'emplois de seniors, moins de pensions. \\
Rendement & Faible à court terme (0,8\md{} en 2030), croissant ensuite. \\
Effets & PIB +0,70~\% en 2037, +1,3~\% en 2047 ; chômage +0,3 point en 2030 le temps de l'absorption. \\
Risques & Pénibilité, inégalités d'espérance de vie (non représentées) ; absorption par le marché du travail plus lente que prévu. \\
Robustesse & \fiab{robuste} pour les pensions ; \fiab{fragile} pour le PIB et l'emploi (fermeture d'offre). \\
\end{longtable}
\end{small}

\paragraph{Ce qui n'est pas recommandé, et pourquoi.}
\begin{itemize}
\item \textbf{Durée hebdomadaire à 37 heures} : effet positif à long terme dans le modèle, mais 400\,000 emplois détruits
à court terme et hypothèse de compensation non estimée.
\item \textbf{Droits de douane} : réduisent la production et le pouvoir d'achat (riposte, prix) et annulent les effets
d'une relocalisation.
\item \textbf{Hausse de l'IS ou des impôts de production} : le modèle n'en voit pas le coût principal (localisation).
\item \textbf{Baisse des allocations chômage} : son effet favorable repose entièrement sur une hypothèse non estimée.
\item \textbf{Retour à 62 ans} : PIB $-$2,5~\% en 2047 ; les pensions supplémentaires ne sont pas financées.
\item \textbf{Réductions d'emplois publics} : coût en activité élevé par euro économisé et effets sur les services non mesurés.
\end{itemize}

\section{Synthèse chiffrée du scénario recommandé à 5, 10 et 20 ans}

\begin{table}[htbp]
\centering\scriptsize
\caption{Scénario R — programme recommandé (modèle complet).}
\label{tab:scR}
\tabfit{../tableaux/scenario_R}
\end{table}

\begin{resultat}[Scénario R : ce que montre le modèle]
\begin{itemize}
\item \textbf{À 5 ans (2032)} : PIB inchangé (+0,02~\%), 63\,000 emplois en moins, chômage +0,5 point (7,5~\%), solde
amélioré de 1,05 point de PIB (environ 31\md), dette plus basse de 3,2 points (138~\% au lieu de 142~\%) ; pouvoir
d'achat $-$1,3~\% (TVA) ; pauvreté $-$0,9 point, indice de Gini $-$1,5 point (statique).
\item \textbf{À 10 ans (2037)} : PIB +0,74~\%, emploi +133\,000, solde +1,3 point, dette $-$7,3 points (171~\% au lieu de
179~\%) ; le chômage reste plus élevé de 0,9 point, le temps que l'économie absorbe les actifs supplémentaires.
\item \textbf{À 20 ans (2047)} : PIB +1,4~\%, emploi +440\,000, solde +2,9 points, dette $-$22 points.
\item \textbf{Distribution} : gagnants, les déciles 1 à 5 (+8,3~\% pour le premier, +0,7~\% pour le cinquième) ; perdants,
les déciles 6 à 10 ($-$0,2 à $-$4,2~\%, le dernier étant le plus touché par l'ISF, le PFU et les droits de succession).
\item \textbf{Robustesse} : les écarts à la référence varient de moins de 0,03 point de PIB et de 0,3 point de dette en
2037 sous les sept chocs ; mis en œuvre à moitié, le programme donne la moitié des gains. Une crise énergétique ou une
hausse des taux de 1,5 point ajoute environ 12 points de dette en 2037, davantage que le gain du programme.
\item \textbf{Limite principale} : \textbf{le programme ne stabilise pas la dette} (171~\% en 2037 dans la référence à taux
constant de 4,89~\%). Il constitue un socle de mesures à faible coût en activité et socialement compensées ; la
stabilisation exigerait, en plus, soit un effort supplémentaire de l'ordre de 4 à 5 points de PIB (dont le coût, dans le
modèle, est une récession de type scénario D), soit une baisse des taux d'intérêt ou une accélération de la productivité
que le modèle ne sait pas produire par lui-même.
\end{itemize}
\end{resultat}

\section{Ce qu'il faudrait savoir de plus}

Les conclusions les plus importantes pour la décision sont aussi celles que le modèle établit le moins bien. Par ordre de
priorité, les améliorations qui changeraient le plus les résultats sont : (1) l'estimation de la fermeture d'offre et du
lien entre production potentielle, travail et productivité (paradoxe des chocs d'offre) ; (2) la réaction de l'offre de
travail aux prestations et aux carrières longues ; (3) un modèle des comportements des entreprises (localisation, marges,
investissement) face à la fiscalité ; (4) les effets de long terme de l'éducation, de la recherche et de la santé sur la
productivité ; (5) des bilans financiers et une prime de risque non linéaire ; (6) une microsimulation dynamique,
revalorisée par le modèle macroéconomique et représentant mieux le haut de la distribution.
